\documentclass[11pt,a4paper]{article}
\usepackage[utf8]{inputenc}
\usepackage[T1]{fontenc}
\usepackage{amsfonts}
\usepackage[french]{babel}
\frenchbsetup{StandardLists=true}
\usepackage{enumitem}
\usepackage{amssymb}
\usepackage{amsmath}
\usepackage[left=2cm,right=2cm,top=2cm,bottom=2cm]{geometry}
\usepackage[dvipsnames]{xcolor}
\usepackage{fouriernc}
\usepackage{graphicx}
\usepackage{setspace}
\singlespacing
\usepackage{tikz}
\usepackage[version=4]{mhchem}
\usepackage{multirow,tabularx,booktabs}
\usepackage{fancyhdr}
\pagestyle{fancy}
\renewcommand\headrulewidth{1pt}
\fancyhead[L]{Lycée Denis Diderot}
\fancyhead[R]{Classe de 1 Bac Pro}
\fancyfoot[C]{page \thepage}
\fancyfoot[R]{Année 2024-2025}
\fancyfoot[L]{Alexis RUIZ}

\renewcommand{\arraystretch}{3.5}

\begin{document}

\begin{center}
\LARGE{Évaluation -- Masses molaires}\\[8mm]
\end{center}
\normalsize
\hrule

\flushleft \textbf{NOM : \hfill Classe : \hspace{3cm} ~~\\[2mm]
Prénom :\\[0,4cm]}
\hrule\vspace{0.2cm}

\vspace{8mm}

\fbox{\textbf{Calcul de masses molaires}}\par\vspace{4mm}

Calculer les masses molaires des composés suivants (on utilisera les données ci-après) :

\begin{enumerate}
\item Éthanol : \ce{C_2 H_6 O}
\item Acide acétique : \ce{C_2 H_4 O_2}
\item Caféine : \ce{C_8 H_{10} N_4 O_2}
\item Paracétamol : \ce{C_8 H_9 N O_2}
\item Saccharose : \ce{C_{12} H_{22} O_{11}}
\end{enumerate}


\begin{center}
\begin{tikzpicture}
\draw (0,0)--(0,3)--(2,3)--(2,0)--cycle;
\node at (1,1.75) {{\LARGE H}};
\node at (1,0.3) {1,0};
\node at (1,1) {Hydrogène};
\node at (0.3,2.6){{\Large 1}};
\end{tikzpicture}
\hspace{0.8cm}
\begin{tikzpicture}
\draw (0,0)--(0,3)--(2,3)--(2,0)--cycle;
\node at (1,1.75) {{\LARGE C}};
\node at (1,0.3) {12,0};
\node at (1,1) {Carbone};
\node at (0.3,2.6){{\Large 6}};
\end{tikzpicture}
\hspace{0.8cm}
\begin{tikzpicture}
\draw (0,0)--(0,3)--(2,3)--(2,0)--cycle;
\node at (1,1.75) {{\LARGE N}};
\node at (1,0.3) {14,0};
\node at (1,1) {Azote};
\node at (0.3,2.6){{\Large 7}};
\end{tikzpicture}
\hspace{0.8cm}
\begin{tikzpicture}
\draw (0,0)--(0,3)--(2,3)--(2,0)--cycle;
\node at (1,1.75) {{\LARGE O}};
\node at (1,0.3) {16,0};
\node at (1,1) {Oxygène};
\node at (0.3,2.6){{\Large 8}};
\end{tikzpicture}
\end{center}

\vspace{5mm}

\noindent\textbf{1.} \dotfill

\vspace{8mm}
\noindent\dotfill

\vspace{8mm}

\noindent\textbf{2.} \dotfill

\vspace{8mm}
\noindent\dotfill

\vspace{8mm}

\noindent\textbf{3.} \dotfill

\vspace{8mm}
\noindent\dotfill

\vspace{8mm}

\noindent\textbf{4.} \dotfill

\vspace{8mm}
\noindent\dotfill

\vspace{8mm}

\noindent\textbf{5.} \dotfill

\vspace{8mm}
\noindent\dotfill

\vspace{8mm}

%\begin{correction}
% 1. Éthanol C2H6O     : M = 2×12,0 + 6×1,0 + 16,0 = 46,0 g/mol
% 2. Acide acétique C2H4O2 : M = 2×12,0 + 4×1,0 + 2×16,0 = 60,0 g/mol
% 3. Caféine C8H10N4O2 : M = 8×12,0 + 10×1,0 + 4×14,0 + 2×16,0 = 194,0 g/mol
% 4. Paracétamol C8H9NO2 : M = 8×12,0 + 9×1,0 + 14,0 + 2×16,0 = 151,0 g/mol
% 5. Saccharose C12H22O11 : M = 12×12,0 + 22×1,0 + 11×16,0 = 342,0 g/mol
%\end{correction}

\end{document}
